\documentclass[a4paper]{article}
% Español (México) con XeLaTeX: fontspec para fuentes Unicode, polyglossia
% para la división silábica y los nombres del documento en la variante mexicana.
\usepackage{fontspec}
\usepackage{polyglossia}
\setdefaultlanguage[variant=mexican]{spanish}
% Los nombres chinos van como «pinyin (original)»: xeCJK da a los caracteres
% Han su propia fuente; sin ella XeLaTeX los omite con un aviso.
% FandolSong no cubre algunos caracteres de nombres (U+73FA); WenQuanYi Zen Hei sí.
\usepackage[AutoFallBack=true]{xeCJK}
\setCJKmainfont{FandolSong-Regular.otf}
\setCJKfallbackfamilyfont{\CJKrmdefault}{WenQuanYi Zen Hei}
% polyglossia no traduce los nombres de listings.
\AtBeginDocument{\providecommand{\lstlistingname}{}\renewcommand{\lstlistingname}{Listado}%
  \providecommand{\lstlistlistingname}{}\renewcommand{\lstlistlistingname}{Índice de listados}}
\usepackage{amsmath, amssymb}
\usepackage{graphicx}
\usepackage[margin=2.35cm]{geometry}
\usepackage[most]{tcolorbox}
\usepackage{booktabs}
\usepackage{tabularx}
\usepackage{longtable}
\usepackage{array}
\usepackage{float}
\usepackage{xcolor}
\usepackage{hyperref}
\usepackage{fancyhdr}
\setCJKmainfont[AutoFakeBold=2.4]{AR PL UMing CN}
\setCJKsansfont[AutoFakeBold=2.4]{Droid Sans Fallback}
\setCJKmonofont[AutoFakeBold=2.4]{Droid Sans Fallback}
\hypersetup{colorlinks=true, linkcolor=blue!55!black, urlcolor=blue!60!black}
\graphicspath{{./}{figures/}}
\setlength{\parskip}{0.55em}
\setlength{\parindent}{2em}
\setlength{\emergencystretch}{3em}
\renewcommand{\arraystretch}{1.18}
\newtcolorbox{knowledgebox}[1]{enhanced, colback=blue!5!white, colframe=blue!70!black, colbacktitle=blue!70!black, coltitle=white, fonttitle=\bfseries, title={#1}, sharp corners, breakable}
\newtcolorbox{importantbox}[1]{enhanced, colback=yellow!10!white, colframe=yellow!75!black, colbacktitle=yellow!75!black, coltitle=black, fonttitle=\bfseries, title={#1}, sharp corners, breakable}
\newtcolorbox{warningbox}[1]{enhanced, colback=red!5!white, colframe=red!70!black, colbacktitle=red!70!black, coltitle=white, fonttitle=\bfseries, title={#1}, sharp corners, breakable}
\pagestyle{fancy}
\fancyhf{}
\fancyfoot[C]{\thepage}
\fancyfoot[R]{\footnotesize\color{gray}\href{https://github.com/hqhq1025/ai-course-notes}{AI Course Notes}}
\renewcommand{\headrulewidth}{0pt}
\renewcommand{\footrulewidth}{0pt}
\newcommand{\videourl}{https://www.youtube.com/watch?v=8dKBH4x0D9o}
\newcommand{\repourl}{https://github.com/hqhq1025/ai-course-notes}

\begin{document}
\begin{titlepage}
\centering
{\Large Notas didácticas de revisión técnica\par}
\vspace{1.0cm}
{\huge\bfseries Repaso de los informes técnicos de Agent: Kimi K2, ChatGPT Agent, Qwen3-Coder y Manus}\par
\vspace{0.5cm}
{\Large Zhang Xiaojun conversa con Zheng Boyuan (郑博元)\par}
\vspace{0.35cm}
{\large Elaborado a partir del video público del Zhang Xiaojun Podcast, una transcripción hecha en GPU, los informes técnicos y entradas de blog oficiales, y diagramas conceptuales propios\par}
\vspace{0.3cm}
{\large 2025-07-30\par}
\vspace{0.85cm}
\includegraphics[width=0.88\textwidth,height=0.42\textheight,keepaspectratio]{cover.jpg}\par
\vfill
\begin{tcolorbox}[width=0.92\textwidth, colback=black!2!white, colframe=black!55, sharp corners]
\textbf{Título del video}: 110. Explicación sección por sección del informe de Kimi K2, comparado con ChatGPT Agent, Qwen3-Coder y otros: «El poder de la ingeniería de sistemas»\par
\textbf{Canal del video}: Zhang Xiaojun Podcast\par
\textbf{Duración del video}: 02:20:45\par
\textbf{Enlace del video}: \href{\videourl}{\nolinkurl{\videourl}}\par
\textbf{Descripción de los materiales}: El video no tiene subtítulos disponibles en YouTube; el cuerpo de la nota se elaboró a partir de la transcripción de faster-whisper large-v3 con CUDA, los capítulos del video, la descripción de YouTube, el informe técnico de Kimi K2, la entrada de blog de Qwen3-Coder, la entrada de blog de Manus sobre ingeniería de contexto y diagramas didácticos propios. La imagen fija de portada del video no se repite en el cuerpo de la nota.\par
\textbf{Página del proyecto}: \href{\repourl}{\nolinkurl{\repourl}}\par
\end{tcolorbox}
\end{titlepage}

\tableofcontents
\newpage

\section{Introducción: por qué estos informes apuntan a los Agents}

Esta sección establece primero los objetivos de aprendizaje de todo el episodio. El programa revisa varios de los materiales técnicos sobre Agents más dignos de una lectura atenta de los últimos meses: el informe técnico de Kimi K2, el anuncio de lanzamiento de ChatGPT Agent, la entrada técnica del blog de Qwen3-Coder y la experiencia de Manus en ingeniería de contexto. Aunque parecen venir de empresas y productos distintos, todos apuntan a un mismo tema: los modelos grandes están pasando de «responder preguntas» a «completar tareas dentro de un entorno».

Lo esencial de un Agent no es solo llamar herramientas, sino el ciclo completo: percibir el entorno, decidir, actuar, leer la retroalimentación y corregir la trayectoria. Este ciclo trae consigo todo un conjunto de problemas de ingeniería: datos sintéticos, generación de trayectorias, aprendizaje por refuerzo, recompensas verificables, entornos aislados de seguridad, ingeniería de contexto, KV cache, selección de herramientas, memoria en el sistema de archivos y organización de múltiples Agents.

\begin{figure}[H]
\centering
\includegraphics[width=0.96\textwidth]{report-comparison.png}
\caption{Comparación de los cuatro informes: Kimi, ChatGPT Agent, Qwen3-Coder y Manus enfatizan cada uno una capa distinta. Diagrama conceptual de elaboración propia, organizado a partir de la estructura del video y de los materiales oficiales.}
\end{figure}

\begin{knowledgebox}{Lectura de la figura: cada material aborda una capa}
Kimi K2 aborda el modelo abierto y el agentic training; ChatGPT Agent, el sistema de producto unificado; Qwen3-Coder, el RL y el uso de herramientas en entornos de código; Manus, la ingeniería de contexto de los Agents en producción. En conjunto forman la cadena completa de los Agents: del modelo y el entrenamiento al producto y el despliegue de ingeniería.
\end{knowledgebox}

\begin{importantbox}{Tesis central del episodio}
El progreso de los Agents no consiste en mejorar la capacidad de un solo modelo, sino en ingeniería de sistemas: solo optimizando en conjunto el modelo, el entorno, las tareas, las herramientas, las recompensas, la seguridad y el contexto es posible llevar la inteligencia de producir texto a completar tareas reales.
\end{importantbox}

\begin{knowledgebox}{Orden de estudio de los cuatro materiales}
\begin{tabularx}{\textwidth}{p{0.20\textwidth}p{0.34\textwidth}X}
\toprule
Material & Qué revisar primero & Valor didáctico \\
\midrule
Kimi K2 & MoE, MuonClip, agentic data, joint RL & Ver cómo se entrena un modelo abierto para volverlo Agentic. \\
ChatGPT Agent & Operator + Deep Research + computadora con herramientas & Ver cómo un Agent se convierte en un producto para usuarios. \\
Qwen3-Coder & Code RL, long-horizon RL, Qwen Code & Ver por qué los entornos de código son adecuados para el RL de Agents. \\
Manus & KV cache, enmascaramiento de herramientas, sistema de archivos, conservación de errores & Ver cómo un Agent en producción logra funcionar de forma estable gracias a la ingeniería de contexto. \\
\bottomrule
\end{tabularx}
\end{knowledgebox}

\begin{knowledgebox}{Resumen de cifras clave}
\begin{tabularx}{\textwidth}{p{0.22\textwidth}p{0.34\textwidth}X}
\toprule
Material & Cifras o hechos clave & Implicación didáctica \\
\midrule
Kimi K2 & Cerca de 1T de parámetros totales, 32B de parámetros activos, 15.5T tokens & Los modelos MoE abiertos también entraron al entrenamiento de sistemas a escala muy grande. \\
Kimi K2 & SWE-Bench Verified 65.8, Tau2-Bench 66.1 & Los Agentic benchmarks se vuelven indicadores centrales de los informes de modelos. \\
Qwen3-Coder & 480B MoE, 35B active, 256K/1M context & Un Agent de código necesita un contexto muy largo y comprensión a nivel de repositorio. \\
Qwen3-Coder & 20,000 entornos paralelos para long-horizon RL & El cuello de botella del RL de Agents es escalar los entornos. \\
Manus & Proporción de tokens de entrada/salida de aproximadamente 100:1 & La KV cache es crucial para el costo y la latencia. \\
\bottomrule
\end{tabularx}
\end{knowledgebox}

\subsection{Resumen de la sección}

EP110 es la clase sobre informes técnicos de Agents. Conecta la teoría de Agents de EP115, la perspectiva de empresa de modelos sobre K2 de EP113 y el Agentic Model de nivel empresarial de EP116, y muestra «el poder de la ingeniería de sistemas».
\section{Definición y clasificación de los Agents: percepción + acción}

La introducción mostró que los cuatro materiales apuntan a los Agents; esta sección primero asienta la base conceptual y responde una pregunta concreta: qué sistema cuenta como Agent y cuál es solo un chatbot. La descripción del video da una definición concisa: un Agent es un sistema inteligente capaz de interactuar con el entorno, con capacidad de percepción y capacidad de acción. La percepción incluye observar el entorno, leer la retroalimentación e interpretar el contexto; la acción incluye llamar herramientas, generar salidas, controlar interfaces y modificar variables. El ciclo mínimo es «observar → decidir → actuar».

\begin{figure}[H]
\centering
\includegraphics[width=0.96\textwidth]{agent-definition-loop.png}
\caption{Definición mínima de un Agent: percibir el entorno, decidir, actuar y volver a leer la retroalimentación. Diagrama conceptual propio, elaborado a partir de la conversación de 00:02:00--00:14:50.}
\end{figure}

\begin{knowledgebox}{Lectura de la figura: diferencia entre un Agent y un Chatbot}
Un Chatbot principalmente responde a la entrada; un Agent tiene que cambiar el estado del entorno. En cuanto el sistema empieza a llamar herramientas, leer la retroalimentación y ajustar su plan, deja de ser un sistema conversacional y pasa a ser un sistema de agente.
\end{knowledgebox}

\subsection{Cuatro tipos de Agent}

Con la definición mínima establecida, esta subsección los clasifica según el entorno y la forma de producto. El programa divide los Agents en Coding Agent, Search Agent, Tool-use Agent y Computer-use Agent. Esta clasificación no es una clasificación académica rigurosa, sino una clasificación por producto y entorno: código, búsqueda, herramientas e interfaz de computadora corresponden a distintas affordance y formas de verificación.

\begin{figure}[H]
\centering
\includegraphics[width=0.96\textwidth]{agent-types-map.png}
\caption{Mapa de tipos de Agent: Coding, Search, Tool-use y Computer-use corresponden a entornos distintos. Diagrama conceptual propio, elaborado a partir de la conversación de 00:02:00--00:14:50.}
\end{figure}

\begin{knowledgebox}{Términos clave: tipos de Agent}
\begin{tabularx}{\textwidth}{p{0.22\textwidth}p{0.34\textwidth}X}
\toprule
Tipo & Tareas representativas & Dificultades \\
\midrule
Coding Agent & Autocompletado de código, refactorización, depuración, modificaciones de PR & Contexto del repositorio de código, pruebas, modificaciones de largo alcance. \\
Search Agent & Recuperación, consolidación, informes de investigación & Fiabilidad de las fuentes, citas, eliminación de duplicados y synthesis. \\
Tool-use Agent & Llamar API, funciones, bases de datos y herramientas de negocio & Selección de herramientas, generación de parámetros, recuperación ante errores. \\
Computer-use Agent & Operar el navegador, la GUI y flujos entre aplicaciones & Grounding visual, seguimiento del estado, seguridad de permisos. \\
\bottomrule
\end{tabularx}
\end{knowledgebox}

\subsection{Resumen de la sección}

La definición de un Agent puede ser muy sencilla, pero llevarla a la práctica es muy complejo. La diferencia esencial entre los distintos tipos de Agent está en que enfrentan entornos, herramientas y retroalimentación distintos.
\section{Dos rutas técnicas: In-Context y End-to-End}

La sección anterior definió el Agent y los tipos de entorno; esta sección trata de cómo construir un Agent. El programa divide las rutas en In-Context Learning y End-to-End Training. La ruta in-context se apoya en un modelo preentrenado potente, el prompt, los protocolos de herramientas y la ingeniería de contexto; permite iterar rápido y ofrece mucha flexibilidad. La ruta end-to-end incorpora la conducta en los pesos del modelo; la inferencia es más estable, pero el costo de entrenamiento es alto y construir los entornos es complejo.

\begin{figure}[H]
\centering
\includegraphics[width=0.94\textwidth]{agent-training-routes.png}
\caption{Dos rutas de entrenamiento de un Agent: In-context es flexible; End-to-end es estable, pero de alto costo. Diagrama conceptual propio, elaborado a partir de la conversación de 00:14:50--00:30:57.}
\end{figure}

\begin{knowledgebox}{Lectura de la figura: las dos rutas no son excluyentes}
Muchos sistemas reales combinan ambas: usan in-context para montar el producto con rapidez, usan trayectorias sintéticas y RL para entrenar en el modelo las conductas clave, y después usan la ingeniería de contexto para resolver los detalles de producción.
\end{knowledgebox}

\subsection{Los tres componentes de Agent Training}

Toda ruta técnica termina en el entrenamiento y el despliegue; esta subsección desglosa los tres componentes de Agent training. Sus etapas clave son la síntesis de datos, el aprendizaje por refuerzo y la seguridad. La síntesis de datos genera trajectory de alta calidad; el aprendizaje por refuerzo depende de una task clara y de una verifiable reward; la seguridad requiere sandbox, restricciones de conducta y human-in-the-loop.

\begin{figure}[H]
\centering
\includegraphics[width=0.96\textwidth]{training-ingredients.png}
\caption{Los tres componentes de Agent Training: las trayectorias sintéticas, el aprendizaje por refuerzo y las restricciones de seguridad forman juntos el pipeline de entrenamiento. Diagrama conceptual propio, elaborado a partir de la conversación de 00:28:29--00:30:57.}
\end{figure}

\begin{importantbox}{El cambio de paradigma en los datos de un Agent}
Los datos anotados tradicionales suelen ser input-output; los datos de un Agent se parecen más a environment + task + action trajectory + reward. El modelo no solo aprende respuestas, sino cómo actuar dentro de un entorno.
\end{importantbox}

\begin{knowledgebox}{Comparación de las etapas de entrenamiento}
\begin{tabularx}{\textwidth}{p{0.22\textwidth}p{0.34\textwidth}X}
\toprule
Etapa & Producto & Riesgo principal \\
\midrule
Data Synthesis & Trayectorias de acción de alta calidad y demostraciones de llamadas a herramientas & Si las trayectorias no son realistas, el modelo aprende a actuar un papel. \\
RL & Optimizar la política a partir de la retroalimentación del entorno & La reward está mal diseñada o sufre un hack. \\
Safety & Entorno aislado, permisos, confirmación humana, políticas de rechazo & Si es demasiado laxa, es peligrosa; si es demasiado estricta, la tarea no se puede completar. \\
Evaluation & Tasa de éxito de las tareas, costo, capacidad de recuperación & El benchmark no coincide con las tareas reales. \\
\bottomrule
\end{tabularx}
\end{knowledgebox}

\subsection{Resumen de la sección}

Lo central de Agent training es colocar la tarea dentro de un entorno para que el modelo aprenda a actuar mediante trayectorias, recompensas y restricciones de seguridad. No es solo SFT ni solo prompt engineering.
\section{Kimi K2: Open Agentic Intelligence}

Las dos secciones anteriores trataron el marco general de los Agent; esta sección aborda el primer informe técnico, y su pregunta central es cómo un modelo abierto puede entrenarse de forma sistemática para adquirir capacidades Agentic y acercarse a la frontera de los modelos cerrados. Kimi K2 es un modelo MoE del orden de 1T de parámetros, con unos 32B de parámetros activos. El informe propone el optimizador MuonClip, que usa QK-Clip para resolver la inestabilidad del entrenamiento con Muon sin perder la token efficiency; K2 se preentrenó con 15.5T tokens, y el posentrenamiento incluye agentic data synthesis a gran escala y joint RL.

\begin{figure}[H]
\centering
\includegraphics[width=0.96\textwidth]{kimi-k2-agentic-pipeline.png}
\caption{Kimi K2 Agentic Pipeline: MuonClip, 15.5T tokens, agentic data synthesis y joint RL. Diagrama conceptual de elaboración propia, basado en el informe técnico de Kimi K2 y en la conversación de 00:30:57--00:43:50.}
\end{figure}

\begin{knowledgebox}{Lectura de la figura: lo importante de Kimi K2 no es un solo benchmark}
K2 reúne la base MoE, un optimizador estable, datos a gran escala, la síntesis de trayectorias de herramientas y RL. Su capacidad agentic proviene de un pipeline de entrenamiento completo, no solo de un procesamiento posterior.
\end{knowledgebox}

\subsection{MuonClip y QK-Clip}

Para entender K2 no basta con ver los agentic benchmark; también hay que ver cómo logra un entrenamiento estable. Esta subsección explica MuonClip y QK-Clip. Muon tiene la ventaja de la token efficiency, pero el entrenamiento a gran escala puede volverse inestable; QK-Clip controla el riesgo de divergencia del entrenamiento al acotar los attention logits. El informe subraya que MuonClip permitió que K2 se mantuviera estable durante el entrenamiento con 15.5T tokens.

\begin{knowledgebox}{Términos clave: componentes del entrenamiento de Kimi K2}
\begin{tabularx}{\textwidth}{p{0.22\textwidth}p{0.34\textwidth}X}
\toprule
Componente & Función & Riesgo/importancia \\
\midrule
MoE & Aumenta el total de parámetros, pero controla el cómputo activo & El enrutamiento, la carga de los expertos y la estabilidad del entrenamiento son complejos. \\
MuonClip & Mejora la token efficiency y estabiliza el entrenamiento & Debe manejar la explosión de los attention logits. \\
QK-Clip & Acota los logits que producen las proyecciones de query/key & Es un mecanismo de estabilización de intervención mínima. \\
Agentic data synthesis & Genera trayectorias de uso de herramientas y datos de tareas & El filtrado de calidad y el diseño de la recompensa son decisivos. \\
Joint RL & Mejora las capacidades en entornos reales o sintéticos & Los entornos y la reward tienen un costo alto. \\
\bottomrule
\end{tabularx}
\end{knowledgebox}

\subsection{Resumen de la sección}

El valor didáctico de Kimi K2 es que combina el modelo abierto, MoE, la estabilidad del optimizador, las trayectorias sintéticas de herramientas y RL en un agentic training pipeline.
\section{ChatGPT Agent: un puente entre la investigación y la práctica}

Kimi K2 muestra el pipeline de entrenamiento de un modelo; esta sección pasa a los sistemas de producto, y la pregunta central es cómo un Agent ya entrenado entra en el flujo de trabajo del usuario y asume los riesgos de operar en entornos reales. El segundo material es el anuncio de lanzamiento de ChatGPT Agent. OpenAI integra Operator, Deep Research y las capacidades de conversación de ChatGPT en un sistema de agentes unificado. Este sistema completa tareas mediante su propia computadora virtual, un navegador de texto, un navegador visual, una terminal y acceso a API, y al mismo tiempo conserva la capacidad del usuario de confirmar, tomar el control e interrumpir.

\begin{figure}[H]
\centering
\includegraphics[width=0.96\textwidth]{chatgpt-agent-unified-system.png}
\caption{Sistema unificado de ChatGPT Agent: Operator + Deep Research + ChatGPT + computadora con herramientas. Diagrama conceptual de elaboración propia, basado en la página oficial de OpenAI y en la conversación de 00:43:50--01:53:38.}
\end{figure}

\begin{knowledgebox}{Lectura de la figura: el sistema de producto tiene muchas más capas que el modelo}
ChatGPT Agent no es solo un modelo: incluye navegación, búsqueda, ejecución de código, manejo de archivos, conectores, confirmación del usuario y políticas de seguridad. Estas capas determinan si puede manejar flujos de trabajo reales.
\end{knowledgebox}

\begin{knowledgebox}{Pila de herramientas de ChatGPT Agent}
\begin{tabularx}{\textwidth}{p{0.22\textwidth}p{0.34\textwidth}X}
\toprule
Herramienta/capa & Función & Control de riesgos \\
\midrule
Navegador visual & Interactúa con páginas web hechas para personas: hace clic, filtra, llena formularios & El usuario puede tomar el control del navegador. \\
Navegador de texto & Lee con eficiencia grandes cantidades de texto web & Requiere filtrado de fuentes y citas. \\
Terminal & Ejecuta código y procesa archivos y datos & Las operaciones de alto riesgo deben restringirse. \\
Connectors/API & Se conecta a Gmail, GitHub, calendario, etc. & Control de permisos y de privacidad. \\
User confirmation & Confirma antes de ejecutar operaciones importantes & Reduce las consecuencias de una operación errónea. \\
\bottomrule
\end{tabularx}
\end{knowledgebox}

\subsection{Seguridad: nuevas capacidades traen nuevos riesgos}

La página de OpenAI enfatiza que este nuevo tipo de Agent puede ejecutar acciones en la web, por lo que trae problemas de inyección de prompts, operaciones erróneas, privacidad y tareas de alto riesgo. Las medidas de control incluyen confirmación de operaciones importantes, modo de toma de control, modo de supervisión, rechazo de tareas de alto riesgo, restricción del acceso a datos y toma de control segura del navegador.

\begin{warningbox}{El cambio central en la seguridad de los Agents}
Cuando un modelo de chat dice algo incorrecto, las consecuencias suelen quedarse en el nivel del texto; cuando un Agent ejecuta una acción incorrecta, puede afectar cuentas, archivos, correos, compras, código y operaciones reales del negocio. Por ello, la seguridad se amplía de la moderación de contenido a los permisos, las herramientas, la confirmación y la auditoría.
\end{warningbox}

\begin{knowledgebox}{Niveles de riesgo de seguridad}
\begin{tabularx}{\textwidth}{p{0.20\textwidth}p{0.34\textwidth}X}
\toprule
Riesgo & Ejemplo & Enfoque de protección \\
\midrule
Inyección de prompts & Instrucciones ocultas en una página web inducen al Agent a filtrar datos & Detección, aislamiento, confirmación de acciones importantes. \\
Operación errónea & Enviar un correo, comprar o borrar un archivo por error & Confirmación humana, toma de control, niveles de permisos. \\
Fuga de privacidad & Se leen datos de conectores o de sitios con sesión iniciada & Mínimo privilegio, limpieza de sesiones, desactivación de conectores no relacionados. \\
Tareas de alto riesgo & Operaciones financieras, biológicas o peligrosas & Políticas de rechazo, modo de supervisión, protecciones específicas. \\
\bottomrule
\end{tabularx}
\end{knowledgebox}

\subsection{Resumen de la sección}

Lo central de ChatGPT Agent es un sistema de producto unificado: reúne en un mismo flujo de trabajo las capacidades de investigación, la operación en la web, la ejecución de herramientas y el control del usuario, y muestra el camino de los Agents desde el experimento hasta el producto.
\section{Qwen3-Coder: Agentic Coding in the World}

La sección anterior examinó los Agent de producto de propósito general; esta sección examina los Agent en entornos de código. La pregunta central es por qué el coding es uno de los entornos de alto valor más adecuados para entrenar Agent y cómo los modelos de código abierto pueden alcanzar a los demás en este entorno. El tercer material es Qwen3-Coder. La publicación oficial del blog presenta Qwen3-Coder-480B-A35B-Instruct: MoE de 480B, 35B active, 256K de context nativo, ampliable a 1M mediante YaRN. Está orientado al agentic coding y destaca el code RL, el long-horizon RL, 20,000 entornos en paralelo y la cadena de herramientas Qwen Code.

\begin{figure}[H]
\centering
\includegraphics[width=0.96\textwidth]{qwen3-coder-stack.png}
\caption{Qwen3-Coder Stack: MoE de 480B, context de 256K/1M, Code RL y Long-horizon RL. Diagrama conceptual de elaboración propia, a partir de la publicación oficial del blog de Qwen y de la conversación de 01:53:38--01:59:04.}
\end{figure}

\begin{knowledgebox}{Lectura de la figura: el código es uno de los entornos más adecuados para el RL de Agent}
Las tareas de código son hard to solve, easy to verify: resolverlas es difícil, pero las pruebas y la retroalimentación de la ejecución son relativamente claras. Esto convierte al código en un buen entorno para el RL a gran escala y la long-horizon interaction.
\end{knowledgebox}

\subsection{Code RL y Long-horizon RL}

La subsección anterior presentó la pila del modelo Qwen3-Coder; esta explica por qué pone el énfasis en el RL. La publicación del blog de Qwen3-Coder subraya que las tareas de código se prestan de forma natural al execution-driven RL. Las tareas reales de ingeniería de software requieren interacción de varios turnos: planificar, usar herramientas, recibir retroalimentación y corregir decisiones. Qwen3-Coder introduce el long-horizon RL y construye entornos a gran escala que pueden ejecutarse en paralelo.

\begin{importantbox}{Por qué el coding es una posición estratégica para los Agent}
Los entornos de código ofrecen acciones claras, retroalimentación ejecutable, verificación mediante pruebas y tareas de alto valor. Frente a muchas tareas de mundo abierto, en el coding es más fácil definir la reward y también es más fácil escalar el entrenamiento y la evaluación.
\end{importantbox}

\begin{knowledgebox}{Por qué el Code RL es «difícil de resolver pero fácil de verificar»}
En muchas tareas de código el proceso de solución es difícil: exige entender la base de código, planificar los cambios y depurar; pero el resultado puede verificarse mediante pruebas, compilación, lint, benchmark o scripts del usuario. Esta estructura es muy adecuada para el aprendizaje por refuerzo, porque la reward puede automatizarse en buena medida.
\end{knowledgebox}

\subsection{Resumen de la sección}

Qwen3-Coder muestra la ruta sistemática de los modelos de código abierto en el agentic coding: contexto largo, datos de código, RL impulsado por la ejecución, interacción con herramientas a largo plazo y una cadena de herramientas de CLI.
\section{Manus: la ingeniería de contexto es el núcleo de un Agent en producción}

Kimi, ChatGPT Agent y Qwen muestran, respectivamente, RL aplicado al modelo, al producto y al código; esta sección examina la capa más subestimada de un Agent en producción: la ingeniería de contexto. Manus optó por construir su Agent sobre la capacidad in-context de los modelos de frontera, en lugar de entrenar desde cero un modelo de extremo a extremo. Subraya que la ingeniería de contexto es una ciencia experimental y que lo decisivo en un Agent en producción no es solo el prompt, sino el KV cache, la gestión de herramientas, la memoria en el sistema de archivos, la reiteración del todo, la conservación de los errores y evitar la trampa de los pocos ejemplos.

\begin{figure}[H]
\centering
\includegraphics[width=0.96\textwidth]{manus-context-engineering.png}
\caption{Manus Context Engineering: KV cache, enmascaramiento de herramientas, sistema de archivos, todo, conservación de errores y diversidad. Diagrama conceptual propio, elaborado a partir de la publicación oficial del blog de Manus y de la conversación de 01:59:04--02:06:06.}
\end{figure}

\begin{knowledgebox}{Lectura de la figura: la ingeniería de contexto es ingeniería en tiempo de ejecución}
El entrenamiento le da capacidades al modelo; la ingeniería de contexto permite que el Agent las use en producción de forma estable, económica y recuperable. Se ocupa de cada ronda de llamadas a herramientas, del crecimiento del contexto y de la recuperación ante fallos.
\end{knowledgebox}

\subsection{KV cache y enmascaramiento de herramientas}

La sección anterior planteó que la ingeniería de contexto es ingeniería en tiempo de ejecución; esta desglosa los dos principios más prácticos: el KV cache y el enmascaramiento de herramientas. Manus considera que la tasa de aciertos del KV-cache es uno de los indicadores más importantes de un Agent en producción, porque la proporción entre entrada y salida de un Agent está muy cargada hacia la entrada, y la caché afecta directamente la latencia y el costo. También propone «enmascarar, no retirar» herramientas: agregar y quitar herramientas de forma dinámica rompe la caché y confunde al modelo; es mejor conservar las definiciones de las herramientas y controlar cuáles pueden elegirse mediante una máquina de estados o un logits mask.

\begin{knowledgebox}{Términos clave: puntos centrales de la ingeniería de contexto}
\begin{tabularx}{\textwidth}{p{0.22\textwidth}p{0.34\textwidth}X}
\toprule
Técnica & Función & Práctica errónea \\
\midrule
KV cache & Reduce el costo de repetir la inferencia sobre prefijos largos & Cambiar el prompt de sistema o las definiciones de herramientas en cada ronda. \\
Tool masking & Controla las acciones disponibles en el momento & Eliminar herramientas de forma dinámica, lo que vuelve inconsistente el contexto. \\
Filesystem as context & Usa el sistema de archivos como memoria externa de largo plazo & Meter todo el contenido en la ventana de contexto. \\
Todo.md & Dirige la atención mediante la reiteración de los objetivos & No actualizar los objetivos en tareas largas. \\
Keep errors & Conserva la evidencia de los fallos para facilitar la recuperación & Limpiar los errores, lo que lleva a repetirlos. \\
\bottomrule
\end{tabularx}
\end{knowledgebox}

\begin{importantbox}{Lista de verificación de un Agent en producción}
Un Agent en producción debe responder al menos seis preguntas: ¿el contexto puede almacenarse en caché? ¿El espacio de herramientas es controlable? ¿Dónde reside la memoria de largo plazo? ¿Se conservan los fallos? ¿Los objetivos se reiteran en el contexto reciente? ¿Los ejemplos repetidos pueden encasillar al modelo en un patrón? Casi toda la experiencia de Manus gira en torno a estas preguntas.
\end{importantbox}

\begin{knowledgebox}{Checklist del sistema de un Agent en producción}
\begin{tabularx}{\textwidth}{p{0.22\textwidth}p{0.34\textwidth}X}
\toprule
Elemento & Pregunta & Consecuencia de omitirlo \\
\midrule
Caché & ¿El prefijo es estable y el contexto es append-only? & El costo y la latencia se disparan. \\
Herramientas & ¿Las definiciones de herramientas son estables y las acciones pueden restringirse? & Elección de la herramienta equivocada, invalidación de la caché, alucinación de acciones. \\
Memoria & ¿El estado de largo plazo se externaliza en archivos o en una base de datos? & Explosión del contexto o pérdida de información. \\
Errores & ¿Las trayectorias fallidas se conservan para que el modelo aprenda de ellas? & Se repite el mismo error. \\
Seguridad & ¿Las acciones de alto riesgo se confirman y se auditan? & Una operación errónea del Agent causa daños reales. \\
\bottomrule
\end{tabularx}
\end{knowledgebox}

\subsection{Resumen de la sección}

Manus nos recuerda que los problemas de un Agent en producción muchas veces no se deben a que el modelo carezca de capacidad, sino a que el contexto, las herramientas, la memoria y la recuperación ante fallos no están bien diseñados.
\section{Nuevo paradigma: Environment + Task-Reward}

Esta sección cierra con la perspectiva planteada en el programa: el núcleo de los datos podría pasar de las anotaciones input-output a la construcción de environment y task-reward. Es decir, en el futuro el entrenamiento de un Agent no consistirá solo en recopilar respuestas, sino en construir entornos interactivos, definir tareas y diseñar reward verificables para que el modelo haga self-improve a partir de la experiencia.

\begin{figure}[H]
\centering
\includegraphics[width=0.96\textwidth]{new-agent-paradigm.png}
\caption{Nuevo paradigma de Agent: el núcleo de los datos pasa de input-output a environment + task-reward. Diagrama conceptual propio, elaborado a partir de la conversación de 02:06:06--02:15:20.}
\end{figure}

\begin{knowledgebox}{Lectura de la figura: de datos de respuestas a datos de experiencia}
Los datos input-output enseñan al modelo a «ver una entrada y dar una respuesta»; environment + task-reward le enseñan a «actuar en un entorno y mejorar a partir de la retroalimentación». Esta es la diferencia de paradigma de datos entre un Agent y un modelo conversacional común.
\end{knowledgebox}

\subsection{Rubrics as reward y self-improvement}

La sección anterior trató el nuevo paradigma de datos; esta se centra en su problema central: de dónde sale la reward. El programa menciona rubrics as reward, es decir, usar criterios de evaluación como mecanismo de recompensa. Su importancia radica en que muchas tareas reales no tienen un simple correcto o incorrecto, pero sí pueden tener criterios de evaluación por niveles. Que un Agent pueda mejorarse a sí mismo depende de que logre encontrar o construir una reward verificable y aprovechar de forma eficaz la experiencia de interacción.

\begin{warningbox}{Condiciones previas de la self-improve}
La automejora de un Agent no significa dejar que el modelo se autocomplazca sin límite. Requiere tareas verificables, entornos confiables, trayectorias que admitan auditoría y mecanismos de seguridad que prevengan el reward hacking.
\end{warningbox}

\begin{knowledgebox}{Tabla comparativa de los nuevos paradigmas}
\begin{tabularx}{\textwidth}{p{0.22\textwidth}p{0.34\textwidth}X}
\toprule
Paradigma & Forma de los datos & Cuello de botella principal \\
\midrule
Input-output & Entradas y respuestas de referencia & Costo de anotación y límites de generalización. \\
Trajectory & Observaciones, acciones, llamadas a herramientas, resultados & Calidad de las trayectorias y recuperación de errores. \\
Environment & Mundo de tareas interactivo & Costo de construcción del entorno y cobertura. \\
Task-reward & Definición de tareas y retroalimentación verificable & Diseño de la reward y seguridad. \\
Experience reuse & Automejora a partir de la experiencia histórica & Memoria, eliminación de ruido y prevención de la autocomplacencia. \\
\bottomrule
\end{tabularx}
\end{knowledgebox}

\subsection{Family of Agents}

La reward y la self-improvement tratan de cómo aprende un sistema individual; esta sección pasa a la organización de múltiples Agent. Al final se menciona que el Agent es como un «cerebro extendido» respaldado por un ejército. La metáfora indica que el Agent del futuro quizá no sea una sola entidad, sino un conjunto de agentes especializados: coding, search, browser, memory, safety, entre otros, que colaboran entre sí.

\begin{figure}[H]
\centering
\includegraphics[width=0.92\textwidth]{family-of-agents.png}
\caption{Family of Agents: el Agent es como un cerebro extendido, respaldado por un ejército. Diagrama conceptual propio, elaborado a partir de la conversación de 02:15:20--02:16:41.}
\end{figure}

\subsection{Resumen de la sección}

El nuevo paradigma de Agent integra datos, entornos, recompensas y ciclos de experiencia. La verdadera ingeniería de sistemas consiste en lograr que estos ciclos sean escalables, verificables y desplegables de forma segura.
\section{Síntesis y ampliación}

Esta sección condensa el EP110 en unas cuantas conclusiones. Primera: la definición mínima de un Agent es percepción y acción, no conversación. Segunda: Kimi K2, ChatGPT Agent, Qwen3-Coder y Manus muestran, respectivamente, las capas de modelo, producto, código/RL e ingeniería de contexto de la pila de un Agent. Tercera: el núcleo del Agent training pasa de input-output a trajectory, environment, task-reward y safety. Cuarta: lo decisivo en un Agent de producción es la ingeniería de sistemas: KV cache, enmascaramiento de herramientas, memoria en el sistema de archivos, recuperación ante errores y control por parte del usuario.

\begin{importantbox}{Ubicar el EP110 en la serie de IA de Zhang Xiaojun (张小珺)}
El EP115 presenta la teoría de la segunda mitad de los Agent; el EP113 aborda Kimi K2 desde la perspectiva de la empresa que desarrolla el modelo; el EP116 trata el Agentic Model empresarial. El EP110, en cambio, reúne varios informes técnicos recientes y muestra la cadena completa de un Agent, desde el artículo y el modelo hasta el producto y su puesta en práctica de ingeniería.
\end{importantbox}

\begin{knowledgebox}{Relación con los episodios anteriores y posteriores}
\begin{tabularx}{\textwidth}{p{0.18\textwidth}p{0.34\textwidth}X}
\toprule
Episodio & Tema & Conexión con el EP110 \\
\midrule
EP115 & Teoría de la segunda mitad de los Agent & Ofrece el marco teórico de reward, environment e interface. \\
EP113 & Kimi K2 y Agentic LLM & Muestra cómo una empresa de modelos entiende K2 y el ecosistema de código abierto. \\
EP116 & Agentic Model empresarial & Ofrece la perspectiva de despliegue de un Agent con datos privados en el ámbito ToB. \\
EP139 & Historia técnica de los Agent & Traza la evolución histórica de los Agent, desde los primeros sistemas hasta los LLM Agent. \\
\bottomrule
\end{tabularx}
\end{knowledgebox}

\subsection{Conclusiones clave}

Las secciones anteriores ya explicaron el Agent desde la definición, el entrenamiento, el producto, el entorno de código y la ingeniería de contexto. Esta sección condensa ese contenido en unos cuantos juicios de ingeniería, para facilitar su comparación posterior con el EP115, el EP113 y el EP116.

\begin{enumerate}
\item Agent = percepción + acción + retroalimentación; no es una capacidad aislada de «saber invocar herramientas».
\item Lo central de Kimi K2 es el modelo MoE abierto, MuonClip, las trayectorias sintéticas de uso de herramientas y el joint RL.
\item Lo central de ChatGPT Agent es el sistema de producto unificado y el control por parte del usuario.
\item Lo central de Qwen3-Coder es el entorno de código, el Code RL, el Long-horizon RL y la cadena de herramientas.
\item Lo central de Manus es la ingeniería de contexto, en particular el KV cache, el enmascaramiento de herramientas y la memoria en el sistema de archivos.
\end{enumerate}

\subsection{Preguntas abiertas}

Estas preguntas son los aspectos que siguen sin resolverse una vez que el Agent pasa de los informes al producto. De ellas dependen el próximo paradigma de entrenamiento, la arquitectura de producto y la dirección de la inversión en infraestructura.

\begin{enumerate}
\item En el entrenamiento de un Agent, ¿en qué proporción deben combinarse las trayectorias sintéticas y la retroalimentación de entornos reales?
\item ¿Pueden las Rubrics as reward sostener una self-improvement confiable en tareas abiertas?
\item ¿En qué forma convergerán finalmente las dos rutas, la in-context y la end-to-end?
\item En un Agent de producción, ¿el principal cuello de botella de costo vendrá de las llamadas al modelo, de la ejecución de herramientas o de la longitud del contexto?
\end{enumerate}

\subsection{Lecturas adicionales}

\begin{itemize}
\item Si le interesa el marco teórico de los Agent, puede compararse con la entrevista a Yao Shunyu (姚顺雨) del EP115.
\item Si le interesa la perspectiva de Kimi K2 como empresa de modelos, puede compararse con la entrevista a Yang Zhilin (杨植麟) del EP113.
\item Si le interesa el Agentic Model empresarial, puede compararse con la entrevista a Wu Minghui (吴明辉) del EP116.
\end{itemize}

\end{document}
